% !TeX root = ../navier-stokes-manuscript.tex
% ============================================================
% 2.1 Vortex Stretching
% ============================================================

The momentum equation has four pieces:
\[
  \underbrace{\partial_t u}_{\text{local acceleration}}
  +
  \underbrace{(u\cdot\nabla)u}_{\text{nonlinear transport}}
  =
  \underbrace{-\nabla p}_{\text{pressure}}
  +
  \underbrace{\nu\Delta u}_{\text{viscous diffusion}}.
\]
The central regularity problem lies in the competition between nonlinear transport and viscosity.
The hope is simple: viscosity smooths the fluid faster than nonlinear transport can sharpen it.
If that remains true at every scale, the velocity stays smooth.

\subsection{Vortex Stretching}\label{sub:ThePerfectlyStretchingVortex}

A narrow material vortex tube is carried by the flow while the same rotating fluid inside it is pulled apart along the tube axis. The symmetric part of the velocity gradient records that axial deformation:
\[
  S:=\frac12\bigl(\nabla u+(\nabla u)^{\mathsf T}\bigr),
\]
and if \(e\) points along the tube segment while \(L(t)\) is its length, extension along \(e\) at rate \(\sigma(t)>0\) means
\[
  Se=\sigma(t)e,
  \qquad
  \sigma(t)>0,
  \qquad
  \frac{\dot L(t)}{L(t)}
  =
  e\cdot Se
  =
  \sigma(t).
\]
The same material segment lengthens at the fractional rate \(\sigma(t)\).

That gain in length is paid for across the tube. Let \(\mathcal V\) be the fixed material volume of the segment, set \(A(t):=\mathcal V/L(t)\), and write \(r_{\mathrm{eff}}(t):=\sqrt{A(t)/\pi}\). Then incompressibility gives
\[
  \nabla\cdot u=0,
  \qquad
  \frac{d}{dt}(A L)=0,
  \qquad
  \frac{\dot A}{A}=-\sigma(t),
  \qquad
  \frac{\dot r_{\mathrm{eff}}}{r_{\mathrm{eff}}}
  =
  -\frac{\sigma(t)}{2}.
\]
As the axis stretches, the core narrows with it.

The rotation in that same fluid is measured by the vorticity \(\omega:=\nabla\times u\), and along the moving tube it obeys
\[
  \frac{D\omega}{Dt}
  =
  S\omega+\nu\Delta\omega,
  \qquad
  \frac{D}{Dt}
  :=
  \partial_t+u\cdot\nabla,
\]
so when the vorticity is aligned with the stretching direction, \(\omega=|\omega|e\),
\[
  S\omega
  =
  \sigma(t)\omega,
  \qquad
  \frac12\frac{D|\omega|^2}{Dt}
  =
  \sigma(t)|\omega|^2
  +
  \nu\,\omega\cdot\Delta\omega.
\]
Now the same positive strain that lengthens the segment and narrows its core also amplifies its aligned vorticity. Viscosity is already in that balance, but the identity does not force it to dominate pointwise.

If the speed stays bounded by \(U_0\) on the fixed-volume material segment \(\Omega_{\mathrm{seg}}(t)\), then
\[
  E_{\mathrm{seg}}(t)
  :=
  \frac12\int_{\Omega_{\mathrm{seg}}(t)}|u(x,t)|^2\,dx
  \leq
  \frac12U_0^2\mathcal V
  <
  \infty,
  \qquad
  \left|\frac12\bigl(\nabla u-(\nabla u)^{\mathsf T}\bigr)\right|_{\mathrm F}
  =
  \frac{|\omega|}{\sqrt2}
  \leq
  |\nabla u|_{\mathrm F}.
\]
A bounded-speed segment can keep finite kinetic energy while axial strain keeps reducing its radius and the rotating part of the velocity gradient keeps growing in the narrowing core. The gradients rise.

\divider{\NavierStokes{} Scaling}

Fix a time \(t_0\) and write the current tube width as \(\ell:=r_{\mathrm{eff}}(t_0)\). Compare the equation at that width with its standard rescaling; this produces another solution for comparison, not a later state of the same material tube. For \(\lambda>1\), define
\[
  u_\lambda(x,t)
  :=
  \lambda u(\lambda x,\lambda^2t),
  \qquad
  p_\lambda(x,t)
  :=
  \lambda^2p(\lambda x,\lambda^2t).
\]
At time \(\lambda^{-2}t_0\), the corresponding vortex in \(u_\lambda\) has width \(\lambda^{-1}\ell\). Differentiating the rescaled fields gives
\[
\begin{aligned}
  \partial_tu_\lambda(x,t)
  &=
  \lambda^3(\partial_tu)(\lambda x,\lambda^2t),
  \\
  (u_\lambda\cdot\nabla)u_\lambda(x,t)
  &=
  \lambda^3\bigl((u\cdot\nabla)u\bigr)(\lambda x,\lambda^2t),
  \\
  \nabla p_\lambda(x,t)
  &=
  \lambda^3(\nabla p)(\lambda x,\lambda^2t),
  \\
  \nu\Delta u_\lambda(x,t)
  &=
  \lambda^3\nu(\Delta u)(\lambda x,\lambda^2t).
\end{aligned}
\]
Every term picks up the same \(\lambda^3\) factor. Passing to a finer width has changed the scale of the comparison, but it has not tilted the nonlinear--viscous competition toward viscosity.

\divider{Critical Height}

That unchanged balance is still not enough to tell whether the narrowing is becoming more severe, because a norm can move simply because the scale moved. Let \(\Lambda:=(-\Delta)^{1/2}\). At Sobolev height \(s\), a change of variables gives
\[
  \norm{u_\lambda(t)}_{\dot H^s}^2
  =
  \lambda^{2s-1}
  \norm{u(\lambda^2t)}_{\dot H^s}^2.
\]
The factor \(\lambda^{2s-1}\) belongs to the rescaling itself, and it disappears only at \(s=\tfrac12\). Set
\[
  H(t)
  :=
  \frac12\norm{u(t)}_{\dot H^{1/2}}^2
  =
  \frac12\norm{\Lambda^{1/2}u(t)}_2^2.
\]
At that height the scale change contributes nothing of its own. The levels on either side move in opposite directions:
\[
  \norm{u_\lambda(t)}_2^2
  =
  \lambda^{-1}\norm{u(\lambda^2t)}_2^2,
  \qquad
  H_\lambda(t)=H(\lambda^2t),
  \qquad
  \norm{\nabla u_\lambda(t)}_2^2
  =
  \lambda\norm{\nabla u(\lambda^2t)}_2^2.
\]
Kinetic energy falls under concentration and first-derivative energy rises, while the critical height stays fixed. Finite energy still leaves the concentration problem open, and any motion of \(H\) has to come from the equation itself.

\divider{Nonlinear Production versus Viscous Dissipation}

Apply \(\Lambda^{1/2}\) to the momentum equation and pair with \(\Lambda^{1/2}u\). Write the signed nonlinear contribution as
\[
  \mathcal P_{1/2}(t)
  :=
  -\operatorname{Re}
  \pair
    {\Lambda^{1/2}u(t)}
    {\Lambda^{1/2}\bigl((u(t)\cdot\nabla)u(t)\bigr)}_{L^2}.
\]
The pressure pairing vanishes because \(\nabla\cdot u=0\), while the Laplacian contributes \(-\nu\norm{\Lambda^{3/2}u}_2^2\). What remains is the exact law
\[
  H'(t)
  +
  \nu\norm{\Lambda^{3/2}u(t)}_2^2
  =
  \mathcal P_{1/2}(t).
\]
The opening hope now has a scale-neutral form: \(H\) rises only when nonlinear production exceeds viscous dissipation. Under the same rescaling,
\[
  \mathcal P_{1/2}[u_\lambda](t)
  =
  \lambda^2\mathcal P_{1/2}[u](\lambda^2t),
  \qquad
  \nu\norm{\Lambda^{3/2}u_\lambda(t)}_2^2
  =
  \lambda^2\nu\norm{\Lambda^{3/2}u(\lambda^2t)}_2^2.
\]
At critical height they scale in exactly the same way. The finer comparison still gives neither side a relative advantage, and the common \(\lambda^2\) factor is paired with the contracted time coordinate \(t\mapsto\lambda^{-2}t\).

\divider{The Heat Clock}

Viscosity gives that contracted clock a concrete meaning. For the linear heat equation \(v_t=\nu\Delta v\),
\[
  \widehat v(\xi,t+\delta t)
  =
  e^{-\nu|\xi|^2\delta t}\widehat v(\xi,t).
\]
A structure localized to width \(r\) lives at frequencies \(|\xi|\simeq r^{-1}\), so the multiplier changes by order one when \(\nu|\xi|^2\delta t\simeq1\). The corresponding viscous comparison time is
\[
  \tau_\nu(r)
  :=
  \frac{r^2}{\nu}.
\]
This is not the lifetime of the material vortex. It is the time viscosity alone needs to change a structure of width \(r\) by order one. At the finer comparison width,
\[
  \tau_\nu(\lambda^{-1}r)
  =
  \lambda^{-2}\tau_\nu(r),
\]
so the viscous comparison time contracts by the same factor as the time coordinate in the full \NavierStokes{} scaling. The unresolved competition has not changed, but every narrower spatial test comes with a quadratically shorter clock.

\divider{The Zeno Cascade}

Repeat that narrowing through the geometric sequence
\[
  r_n:=2^{-n}r_0,
  \qquad
  \tau_n:=\frac{r_n^2}{\nu},
\]
and the associated clocks satisfy
\[
  \tau_n
  =
  \frac{r_0^2}{\nu}4^{-n},
  \qquad
  \sum_{n=0}^{\infty}\tau_n
  =
  \frac{4r_0^2}{3\nu}
  <
  \infty.
\]
Infinitely many finer comparison scales fit inside finite total comparison time. That is still only a scale--clock schedule. For it to describe a candidate singularity, one actual solution would have to pass through widths
\[
  r_0>r_1>r_2>\cdots\longrightarrow0
\]
at times
\[
  t_0<t_1<t_2<\cdots\uparrow T_*<\infty,
  \qquad
  t_{n+1}-t_n\asymp\tau_n,
\]
with comparison constants independent of \(n\). The remaining time would then satisfy
\[
  T_*-t_n
  \asymp
  \sum_{k=n}^{\infty}\tau_k
  =
  \frac{4r_n^2}{3\nu}.
\]
By stage \(n\), the same velocity history would have only order \(r_n^2/\nu\) of time left before \(T_*\). The spatial schedule says how the widths and clocks could accumulate; it does not say that a smooth solution can make that passage. What it does leave behind is unavoidable: if one history is to cross infinitely many shrinking widths in finite time, the field must keep producing its next response on clocks collapsing at that same quadratic rate.
